\documentclass[11pt]{article}
\usepackage[letterpaper,margin=1in]{geometry}
\usepackage{amsmath,amssymb,amsthm,mathtools}
\usepackage{booktabs,array,longtable}
\usepackage{xcolor}
\usepackage{microtype}
\usepackage[hidelinks]{hyperref}
\hypersetup{
  pdftitle={A frustrated-conditioning obstruction for local Yang--Mills defect domination},
  pdfsubject={Center staples, conditional heat-kernel measures, and dependent polymer strategy}
}

\newtheorem{theorem}{Theorem}[section]
\newtheorem{proposition}[theorem]{Proposition}
\newtheorem{lemma}[theorem]{Lemma}
\newtheorem{corollary}[theorem]{Corollary}
\theoremstyle{definition}
\newtheorem{definition}[theorem]{Definition}
\newtheorem{remark}[theorem]{Remark}

\newcommand{\SU}{\mathrm{SU}(2)}
\newcommand{\Ztwo}{\mathbb Z_2}
\newcommand{\norm}[1]{\left\lVert #1\right\rVert}
\newcommand{\dd}{\,\mathrm{d}}
\newcommand{\proved}{\textcolor{green!45!black}{\textsf{proved}}}
\newcommand{\restricted}{\textcolor{blue!65!black}{\textsf{restricted}}}
\newcommand{\openstatus}{\textcolor{red!70!black}{\textsf{open}}}

\title{A Frustrated-Conditioning Obstruction for Local Yang--Mills Defect Domination\\
\large Why dependent ultraviolet control must charge the boundary to a contour}
\author{Mathematical research dossier}
\date{20 July 2026}

\begin{document}
\maketitle

\begin{abstract}
An ultraviolet one-plaquette heat-kernel tail can be small while a
worst-case conditional tail in the interacting link measure is large.  We
prove this separation exactly on the side-two four-dimensional $\SU$ block.
An interior link belongs to six plaquettes.  The center-staple map from the
remaining links onto those six plaquettes has $\Ztwo$-rank six, so an
admissible center-valued conditioning can give five identity staples and one
negative-identity staple.  The selected link then has conditional density
\[
 \gamma_t(\dd U)\propto p_t(U)^5p_t(-U)\dd U.
\]
Its weak-coupling class-angle rate function is
\[
 I_6(\theta)=5\theta^2+(\pi-\theta)^2.
\]
The unique minimizer is $\theta_*=\pi/6$, leaving the designated frustrated
plaquette at angle $5\pi/6$.  If its good-curvature radius is
$0<\rho<5\pi/6$, then
\[
 \lim_{t\downarrow0}2t\log
 \gamma_t\{\theta_{\mathrm{target}}<\rho\}
 =-6\left(\frac{5\pi}{6}-\rho\right)^2.
\]
Thus the conditional bad probability tends to one.  In particular, no
rare-defect estimate uniform over arbitrary microscopic boundary
conditionings can initialize the Stage 09 independent-polymer benchmark.

This obstruction eliminates a proof mechanism, not the constructive
programme.  A viable dependent estimate must retain the cost of the
frustrating boundary: condition on a good collar, integrate a multi-link
block, or charge the entire connected frustration to a labeled contour that
also records center and topological data.
\end{abstract}

\tableofcontents

\section{The dependence question}

Let $B$ be the cubical complex with vertices $\{0,1,2\}^4$.  It has 216 links
and 216 plaquettes.  With heat generator $\tfrac12\Delta_{\SU}$, the regulated
pure plaquette weight is
\[
 \prod_{p\in B_2}p_t(U_p)\prod_{e\in B_1}\dd U_e.
\]
Stage 09 proved that at $t=3\cdot10^{-6}$ a local nonlinear action chart and a
standalone one-plaquette tail can satisfy the desired inequalities at the same
scale.  The missing step was dependence: could the standalone tail be
replaced by a bound valid under the link interactions?

A tempting sufficient condition would be
\begin{equation}\label{eq:uniform}
 \sup_{\omega}
 \Pr_t\{U_p\text{ leaves the radius-}\rho\text{ chart}\mid\omega\}
 \le\varepsilon(t),\qquad \varepsilon(t)\longrightarrow0,
\end{equation}
where $\omega$ ranges over all surrounding microscopic link conditionings.
Such a bound could feed a Dobrushin or stochastic-domination argument without
tracking the probability of $\omega$.  The central result shows that
\eqref{eq:uniform} has the wrong quantifier.

\begin{theorem}[Frustrated-conditioning obstruction]\label{thm:main}
On the open side-two four-dimensional $\SU$ block, there is an interior link
$e$ and a center-valued assignment of every other link such that five
plaquette staples incident to $e$ equal $I$ and one equals $-I$.  Conditioned
on those other links, $U_e$ has density
\begin{equation}\label{eq:density}
 \gamma_t(\dd U)=Z_t^{-1}p_t(U)^5p_t(-U)\dd U.
\end{equation}
If the plaquette with staple $-I$ is declared good when its class angle is
less than $\rho$, then for every $0<\rho<5\pi/6$,
\begin{equation}\label{eq:limit}
 \lim_{t\downarrow0}2t\log
 \gamma_t\{\theta_{\mathrm{target}}<\rho\}
 =-6\left(\frac{5\pi}{6}-\rho\right)^2<0.
\end{equation}
Consequently
\[
 \lim_{t\downarrow0}
 \sup_{\omega}\Pr_t\{U_p\text{ is bad}\mid\omega\}=1.
\]
\end{theorem}

At the Stage 09 curvature radius $\rho=0.0082$, the rate gap in
\eqref{eq:limit} is $40.8661445$.  The obstruction is therefore separated by
a large action-scale margin, not a borderline numerical constant.

\section{Exact center-staple realization}\label{sec:center}

Choose a link of maximum plaquette degree.  The exact incidence matrix of the
block has six nonzero entries in its column.  Hold the selected link variable
and assign every other link a center sign $(-I)^{x_e}$,
$x_e\in\Ztwo$.  Orientation disappears because $(-I)^{-1}=-I$.  The six
staple signs are therefore a linear map
\[
 C:\Ztwo^{215}\longrightarrow\Ztwo^6
\]
obtained by deleting the selected incidence column and reducing the remaining
six rows modulo two.

\begin{proposition}[Staple surjectivity]\label{prop:rank}
The map $C$ has rank six.  In particular, every prescribed pattern of center
staples on the six incident plaquettes is realized by an assignment of the
remaining links.
\end{proposition}

\begin{proof}
For each incident plaquette, take the link opposite and parallel to the
selected link.  These six opposite links are distinct.  The opposite link for
a given plaquette belongs to that incident plaquette and to none of the other
five incident plaquettes.  The six corresponding columns of $C$ therefore
form the $6\times6$ identity matrix after row ordering.  Hence
$\operatorname{rank}_{\Ztwo}C=6$.
\end{proof}

The executable certificate independently constructs the full $6\times215$
matrix, performs Gaussian elimination over $\Ztwo$, and solves for a vector
with exactly one requested negative staple.  It selects link
\[
 e=((0,1,1,1),0),
\]
whose six incident plaquettes have global indices
$21,44,54,58,59,60$.  A solution uses one negative-center link and verifies
the requested staple vector exactly.

\begin{remark}[Positive-measure stability]
The exact center assignment is a legitimate value of the regular conditional
density.  Moreover, the six staples and the conditional density depend
continuously on the surrounding links.  The displaced conditional minimum and
the bad-probability conclusion persist for a neighborhood of the center
assignment.  The obstruction is therefore not an artifact of selecting a
single measure-zero boundary value.
\end{remark}

\section{The exact one-link conditional law}\label{sec:law}

All plaquette factors not incident to $e$ are constant under conditioning.
For an incident identity staple, conjugation invariance and inversion symmetry
give the factor $p_t(U_e)$.  For the one negative-center staple the factor is
$p_t(-U_e)$.  This proves \eqref{eq:density}.

Write $\theta\in[0,\pi]$ for the class angle of $U_e$.  The class angle of
$-U_e$ is $\pi-\theta$.  Radial Haar integration therefore gives
\begin{equation}\label{eq:radial}
 \gamma_t(\dd\theta)\propto
 \sin^2\theta\,p_t(\theta)^5p_t(\pi-\theta)\dd\theta.
\end{equation}
The target plaquette is good precisely when
\[
 \pi-\theta<\rho,
 \qquad\text{or equivalently}\qquad
 \theta>\pi-\rho.
\]

This conditional law exposes why a one-plaquette marginal is insufficient.
The five aligned plaquettes collectively compete with the designated one.
Weak coupling does not force every factor independently to its own identity;
it minimizes their total action.

\section{Small-time rate theorem}\label{sec:rate}

The $\SU$ heat kernel in the declared normalization has the class-angle
small-time rate
\begin{equation}\label{eq:hk-rate}
 \lim_{t\downarrow0}-2t\log p_t(\theta)=\theta^2,
 \qquad 0\le\theta\le\pi.
\end{equation}
The exact Poisson-image representation proves this directly: the shortest
geodesic image supplies the exponential rate, while all prefactors and the
cut-locus multiplicity contribute only subexponential powers of $t$.

\begin{proposition}[Conditional large-deviation rate]\label{prop:rate}
The probability measures \eqref{eq:radial} obey the class-angle rate function
\begin{equation}\label{eq:I}
 I_6(\theta)=5\theta^2+(\pi-\theta)^2,
\end{equation}
up to subtraction of its minimum.  The unique minimizer is
\[
 \theta_*=\frac\pi6,
 \qquad
 I_6(\theta_*)=\frac{5\pi^2}{6}.
\]
The target plaquette angle at the minimizer is $5\pi/6$.
\end{proposition}

\begin{proof}
Apply \eqref{eq:hk-rate} to the six heat-kernel factors in
\eqref{eq:radial}.  The Haar factor $\sin^2\theta$ is subexponential and does
not alter the rate.  Completing the square gives
\[
 I_6(\theta)
 =6\left(\theta-\frac\pi6\right)^2+\frac{5\pi^2}{6},
\]
which proves every assertion.
\end{proof}

\begin{proof}[Proof of Theorem~\ref{thm:main}]
Proposition~\ref{prop:rank} gives the boundary assignment and
Section~\ref{sec:law} gives the conditional law.  If $\rho<5\pi/6$, the good
region $[\pi-\rho,\pi]$ lies strictly to the right of $\theta_*$.  Its rate
infimum is attained at $\pi-\rho$.  The Laplace principle and
Proposition~\ref{prop:rate} give
\begin{align*}
 &I_6(\pi-\rho)-I_6(\theta_*)\\
 &\qquad=6\left(\pi-\rho-\frac\pi6\right)^2
 =6\left(\frac{5\pi}{6}-\rho\right)^2.
\end{align*}
This proves \eqref{eq:limit}.  The good probability tends to zero, so the bad
probability tends to one.  Taking the supremum over boundary conditionings
can only increase it.
\end{proof}

\section{Executable calibration}

The implementation has two logically distinct parts.  Exact finite algebra
constructs the incidence restriction, computes its $\Ztwo$ rank, solves the
staple equation, and evaluates the closed rate formulas.  Numerical
calibration uses a signed-log Poisson-image sum and logarithmic
Gauss--Legendre quadrature.  It never exponentiates the ultraviolet
probability and therefore remains stable when its logarithm is of order
millions.

\begin{table}[ht]
\centering
\caption{Convergence of the conditional good-event logarithmic rate for
$\rho=0.0082$.}\label{tab:rates}
\begin{tabular}{@{}rrrr@{}}
\toprule
$t$ & mode $\theta$ & $-2t\log P(\mathrm{good})$ & relative error \\
\midrule
$0.05$ & $0.541692$ & $39.578307$ & $3.151\cdot10^{-2}$ \\
$0.02$ & $0.530832$ & $40.274290$ & $1.448\cdot10^{-2}$ \\
$0.01$ & $0.527454$ & $40.570358$ & $7.238\cdot10^{-3}$ \\
$3\cdot10^{-6}$ & $0.523599$ & $40.866026$ & $2.908\cdot10^{-6}$ \\
\bottomrule
\end{tabular}
\end{table}

The exact limiting value is $40.866145$.  At the Stage 09 heat time the
computed log good probability is approximately $-6.811004\cdot10^6$.
Floating-point representation consequently rounds the bad probability to
one, while the log-domain certificate preserves the quantitative separation.
The numerical values calibrate the analytic theorem; they are not used to
prove it.

\section{What proof route has been eliminated?}

The theorem rules out any dependent-defect argument whose decisive input is a
uniform estimate over arbitrary microscopic exterior links, such as
\[
 \sup_{\omega}\Pr(B_p\mid\omega)\le p(t),
 \qquad p(t)\longrightarrow0,
\]
for $B_p=\{\theta_p\ge\rho\}$ with fixed
$\rho<5\pi/6$.  It also rules out inserting the Stage 09 standalone tail into
such an estimate: the standalone law and the frustrated conditional law have
different rate minimizers.

The result does \emph{not} show that bad blocks are frequent under the full
unconditioned measure.  The frustrating boundary may itself have a large
action cost.  Taking a supremum makes that cost disappear, which is precisely
why the method fails.  Nor does the result rule out a conditional theorem
restricted to a good boundary collar.

\section{Replacement: charged and typed contours}\label{sec:replacement}

A viable activity expansion must include the cause of the frustration in the
same connected object whose weight is being estimated.  The minimal typed
defect alphabet should distinguish
\[
 \begin{split}
 \mathsf{Defect}={}&\mathsf{CurvatureExit}(\rho)
 \mid\mathsf{OrbitAtlasFailure}(c)
 \mid\mathsf{CenterFlux}(z)\\
 &\mid\mathsf{BoundaryHolonomy}(h)
 \mid\mathsf{TopologicalCharge}(q)
 \mid\mathsf{InteractionNormFailure}(n).
 \end{split}
\]
The labels cannot be collapsed to a single bad-block bit before their action
and blocking rules have been charged.

The next theorem should combine two estimates with identical parameter
registries.  First, for a gauge-invariant good-collar event $G(A)$ surrounding
a block $B$, prove
\begin{equation}\label{eq:collar}
 \sup_{\omega\text{ compatible with }G(A)}
 \Pr\{B\text{ exits the chart}\mid\omega\}
 \le p_{\mathrm{good}}(t,\rho,A).
\end{equation}
Second, control failure of the collar by labeled activities:
\begin{equation}\label{eq:activity}
 \sup_L\sup_o\sum_{X\ni o}
 e^{\alpha|X|}|z_{\mathrm{collar},L}(X)|<\delta.
\end{equation}
Alternatives that preserve the same information are full multi-link block
integration and reflection-positive chessboard estimates.  In each case the
center-valued frustration must be paid for, not treated as free boundary data.

\section{Frontier audit}

\begin{longtable}{@{}p{0.45\textwidth}p{0.16\textwidth}p{0.29\textwidth}@{}}
\toprule
Claim & Status & Scope \\
\midrule
\endhead
Degree-six interior-link incidence & \proved & finite $2^4$ block \\
$\Ztwo$ center-staple surjectivity & \proved & exact finite algebra \\
Conditional law \eqref{eq:density} & \proved & one admissible conditioning \\
Rate minimizer and gap & \proved & $t\downarrow0$ conditional law \\
Worst-case uniform-tail obstruction & \restricted & microscopic conditionings \\
Good-collar theorem \eqref{eq:collar} & \openstatus & interacting link measure \\
Charged activity sum \eqref{eq:activity} & \openstatus & uniform in volume \\
Repeated RG and basin entry & \openstatus & cutoff and scale uniformity \\
Continuum construction and Hamiltonian gap & \openstatus & Clay endpoint \\
\bottomrule
\end{longtable}

The strongest defensible level is \textbf{C2, restricted obstruction
theorem}.  The result makes the dependent-measure target more precise by
eliminating an over-strong quantifier.  It does not supply the replacement
activity expansion, construct continuum Yang--Mills theory, or prove a mass
gap.

\begin{thebibliography}{9}
\bibitem{MenottiOnofri1981}
P.~Menotti and E.~Onofri,
``The action of $SU(N)$ lattice gauge theory in terms of the heat kernel on the
group manifold,''
\emph{Nuclear Physics B} \textbf{190} (1981), 288--300.

\bibitem{Balaban1987}
T.~Balaban,
``Renormalization group approach to lattice gauge field theories. I.
Generation of effective actions in a small field approximation and a coupling
constant renormalization in four dimensions,''
\emph{Communications in Mathematical Physics} \textbf{109} (1987), 249--301.
\end{thebibliography}

\end{document}

